\documentclass[11pt]{article}
\usepackage[margin=1in]{geometry}
\usepackage{enumitem}
% =============================================================================
% Preambles/header.tex — shared LaTeX/Beamer preamble for this project
%
% Use from a Beamer lecture by prepending:
%     \input{header}
% and compiling with TEXINPUTS=../Preambles:$TEXINPUTS (the /compile-latex
% skill does this automatically).
%
% PALETTE CONTRACT. Color names defined here MUST match the names in
% Quarto/theme-template.scss so Beamer and Quarto renderings use the same
% palette. The scripts/check-palette-sync.sh script enforces this.
%
% To customize: edit the HEX values below AND the matching $variable in
% Quarto/theme-template.scss, then run scripts/check-palette-sync.sh.
% =============================================================================

% -----------------------------------------------------------------------------
% Engine detection + encoding
%
% Our /compile-latex skill uses XeLaTeX, where inputenc is unnecessary and
% can produce encoding warnings. Only load inputenc under pdfTeX.
% -----------------------------------------------------------------------------
\usepackage{iftex}
\ifPDFTeX
  \usepackage[utf8]{inputenc}
\fi

\usepackage{amsmath, amssymb, amsthm}
\usepackage{booktabs}
\usepackage{graphicx}

% -----------------------------------------------------------------------------
% PALETTE — mirrors Quarto/theme-template.scss variable names.
% Load xcolor and define colors BEFORE anything that references them
% (notably hyperref, hypersetup, and the Beamer color assignments).
% -----------------------------------------------------------------------------
\usepackage{xcolor}

\definecolor{primary-blue}{HTML}{012169}      % headings, accents
\definecolor{primary-gold}{HTML}{B9975B}      % emphasis, borders
\definecolor{highlight-yellow}{HTML}{F2A900}  % markers, alerts
\definecolor{light-bg}{HTML}{E8EDF5}          % title / section backgrounds
\definecolor{jet}{HTML}{1A1A1A}               % body text

% Semantic colors (used in diagrams and annotations)
\definecolor{positive}{HTML}{15803D}          % good / observed / identified
\definecolor{negative}{HTML}{B91C1C}          % bad / problematic / bias
\definecolor{neutral}{HTML}{525252}           % reference / context

% Highlight accents (match .hi-* classes in the SCSS)
\definecolor{hi-slate}{HTML}{314F4F}
\definecolor{hi-green}{HTML}{2E8B57}
\definecolor{hi-red}{HTML}{C41E3A}

% Semantic aliases so skill templates and snippets can write
% \textcolor{accent}{...} without hard-coding "primary-blue".
\colorlet{accent}{primary-blue}
\colorlet{accent2}{primary-gold}

% -----------------------------------------------------------------------------
% Hyperref — loaded AFTER xcolor + palette so link colors resolve.
% -----------------------------------------------------------------------------
\usepackage{hyperref}
\hypersetup{colorlinks=true, linkcolor=primary-blue, urlcolor=primary-blue,
            citecolor=primary-blue, pdfborder={0 0 0}}

% -----------------------------------------------------------------------------
% Course code — set once per deck (after \input{header}, before \begin{document})
% via \coursecode{CS401}. Shown in the footer on every slide so a deck's course
% is identifiable without opening the title page. Empty by default.
% -----------------------------------------------------------------------------
\makeatletter
\newcommand{\coursecode}[1]{\gdef\@coursecodetext{#1}}
\gdef\@coursecodetext{}
\makeatother

% -----------------------------------------------------------------------------
% Beamer theme assignments (only apply under Beamer).
%
% We detect Beamer via \@ifclassloaded{beamer}, which is the documented,
% non-internal way. This must be wrapped in \makeatletter/\makeatother
% because \@ifclassloaded contains an @-catcode character.
% -----------------------------------------------------------------------------
\makeatletter
\@ifclassloaded{beamer}{%
  \usefonttheme{professionalfonts}
  \setbeamercolor{structure}{fg=primary-blue}
  \setbeamercolor{frametitle}{fg=primary-blue}
  \setbeamercolor{title}{fg=primary-blue}
  \setbeamercolor{subtitle}{fg=primary-gold}
  \setbeamercolor{author}{fg=primary-blue}
  \setbeamercolor{institute}{fg=neutral}
  \setbeamercolor{date}{fg=primary-gold}
  \setbeamercolor{itemize item}{fg=primary-blue}
  \setbeamercolor{itemize subitem}{fg=primary-gold}
  \setbeamercolor{alerted text}{fg=primary-gold}
  \setbeamercolor{block title}{fg=white, bg=primary-blue}
  \setbeamercolor{block body}{bg=light-bg}
  % Semantic block variants: block=definition/concept (blue), exampleblock=
  % worked example (green/positive), alertblock=key takeaway or caution (gold).
  % Keep to <=2 colored boxes per slide across ALL three variants (INV-7).
  \setbeamercolor{block title example}{fg=white, bg=positive}
  \setbeamercolor{block body example}{bg=light-bg}
  \setbeamercolor{block title alerted}{fg=white, bg=primary-gold}
  \setbeamercolor{block body alerted}{bg=light-bg}

  % Minimal footer: course code (left) + page X / N (right), no navigation clutter
  \setbeamertemplate{navigation symbols}{}
  \setbeamertemplate{footline}{%
    \color{neutral}\footnotesize\hspace{0.7em}\@coursecodetext\hfill
    \insertframenumber/\inserttotalframenumber\hspace{0.7em}\vspace{0.3em}}
}{}
\makeatother

% -----------------------------------------------------------------------------
% TikZ setup — libraries the snippet gallery uses
% -----------------------------------------------------------------------------
\usepackage{tikz}
\usetikzlibrary{arrows.meta, positioning, calc, decorations.pathreplacing,
                fit, shapes.geometric, backgrounds}

% Shared TikZ styles — snippets and hand-written diagrams can pick these up
% via `\begin{tikzpicture}[every node/.style={...}]` or by naming specific
% styles (e.g., `\node[dag-node] (X) at (0,0) {X};`).
\tikzset{
  % Node styles for causal DAGs and similar
  dag-node/.style = {
    draw=primary-blue, fill=light-bg, rounded corners=2pt,
    minimum width=1.2cm, minimum height=0.9cm, align=center, font=\small
  },
  decision-node/.style = {
    draw=primary-gold, fill=white, diamond, aspect=1.8,
    minimum width=1.8cm, minimum height=1.2cm, align=center, font=\small,
    inner sep=1pt
  },
  % Edge styles — observed vs counterfactual
  observed-edge/.style = {-{Latex[length=2.5mm]}, thick, draw=jet},
  counterfactual-edge/.style = {-{Latex[length=2.5mm]}, thick, draw=neutral, dashed},
  confound-edge/.style = {-{Latex[length=2.5mm]}, thick, draw=negative, dashed},
  % Data-point styles for plots
  observed-dot/.style = {fill=primary-blue, draw=primary-blue, circle, inner sep=1.2pt},
  counterfactual-dot/.style = {fill=white, draw=neutral, circle, inner sep=1.2pt}
}

% -----------------------------------------------------------------------------
% Convenience macros
% -----------------------------------------------------------------------------
\newcommand{\muted}[1]{\textcolor{neutral}{#1}}
\newcommand{\key}[1]{\textcolor{primary-gold}{\textbf{#1}}}
\newcommand{\good}[1]{\textcolor{positive}{#1}}
\newcommand{\bad}[1]{\textcolor{negative}{#1}}

% Standout frame macro: full-bleed dark background for section transitions.
% Beamer-only (uses \begin{frame}/\setbeamercolor/\usebeamerfont) -- do not
% call from an article-class file (e.g. Notes/<CODE>/*-notes.tex). \muted,
% \key, \good, \bad, and \coursecode above are plain xcolor/gdef and are
% safe to use from either class; verified when Notes/article-class support
% was added.
% Expand: \transitionslide{Your transition title}
\newcommand{\transitionslide}[1]{%
  \begingroup
  \setbeamercolor{background canvas}{bg=primary-blue}
  \begin{frame}[plain]
    \centering
    \vfill
    {\usebeamerfont{title}\color{white}\Large #1\par}
    \vfill
  \end{frame}
  \endgroup
}

% Section divider frame (Beamer-only), an alternative to \transitionslide with a
% darker two-line style: near-black (jet) background, a small white label line
% above a larger gold title, and a thin gold rule along the bottom edge.
% Expand: \sectiondivider{Section 2}{Pointers}
\newcommand{\sectiondivider}[2]{%
  \begingroup
  \setbeamercolor{background canvas}{bg=jet}
  \begin{frame}[plain]
    \vspace*{3.0cm}
    {\color{white}\ttfamily\large #1\par}
    \vspace{0.5cm}
    {\color{primary-gold}\LARGE #2\par}
    \begin{tikzpicture}[remember picture, overlay]
      \draw[primary-gold, line width=2.5pt]
        ($(current page.south west)+(0,0.1)$) -- ($(current page.south east)+(0,0.1)$);
    \end{tikzpicture}
  \end{frame}
  \endgroup
}

% -----------------------------------------------------------------------------
% Channel watermark — "Concepts That Click" (YouTube).
%
% Article-class files (CompetitiveExam Q/A sets, Notes, Assignments) get a
% light diagonal draft watermark on every page plus a centered footer naming
% the channel. Beamer decks get a subtle TikZ background watermark instead
% (the footline already carries the course code). Centralized here so every
% MCQ PDF is branded without per-file edits.
% -----------------------------------------------------------------------------
\makeatletter
\@ifclassloaded{article}{%
  \usepackage{draftwatermark}
  \SetWatermarkText{Concepts That Click}
  \SetWatermarkScale{1}
  \SetWatermarkFontSize{2cm}
  \SetWatermarkLightness{0.86}
  \SetWatermarkAngle{45}
  \usepackage{fancyhdr}
  \pagestyle{fancy}
  \fancyhf{}
  \fancyfoot[C]{\footnotesize\textcolor{neutral}{Concepts That Click~|~YouTube: \href{https://www.youtube.com/@concepts.thatclick}{@concepts.thatclick}}}
  \renewcommand{\headrulewidth}{0pt}
  \renewcommand{\footrulewidth}{0pt}
  \fancypagestyle{plain}{%
    \fancyhf{}%
    \fancyfoot[C]{\footnotesize\textcolor{neutral}{Concepts That Click~|~YouTube: \href{https://www.youtube.com/@concepts.thatclick}{@concepts.thatclick}}}%
    \renewcommand{\headrulewidth}{0pt}%
    \renewcommand{\footrulewidth}{0pt}%
  }
}{}
\@ifclassloaded{beamer}{%
  \setbeamertemplate{background}{%
    \begin{tikzpicture}[remember picture, overlay]
      \node[rotate=45, opacity=0.055, align=center]
        at (current page.center)
        {\Huge\textcolor{neutral}{Concepts That Click}};
    \end{tikzpicture}%
  }
}{}
\makeatother 
\coursecode{JKSSB}
\renewcommand{\thesection}{09.\arabic{section}}
\newtheorem{example}{Example}[section]
\newenvironment{definitionbox}[1]{%
  \par\noindent\textbf{Definition (#1).}\ \itshape
}{\par\vspace{0.3em}}
\title{Virtual Memory --- Detailed Lecture Notes}
\author{JKSSB Computer Awareness}
\date{\today}

\begin{document}
\maketitle

\section{What virtual memory is and why it exists}
Physical RAM is limited. A modern application, its data, and the other
programs open at the same time can together demand more memory than the
machine physically has. \textbf{Virtual memory} is the technique that lets the
system keep running anyway: it uses secondary storage (the disk) to hold parts
of a program that are not currently needed in RAM. In this way the \emph{addressable}
memory seen by a program can be much larger than the physical memory actually
installed.

\begin{definitionbox}{Virtual memory}
A memory-management technique that uses secondary storage to extend the
addressable memory beyond physical RAM, giving each process the impression of a
large, contiguous memory space.
\end{definitionbox}

The office-desktop running example (EX-01) makes this concrete. The desktop has
16\,GB of RAM and runs MS Word, a spreadsheet, and a web browser at once. If the
combined working set of those programs exceeds 16\,GB, the operating system
moves some of the least-used pages to disk and brings them back when they are
referenced again. No single program is told that RAM ran out; each continues to
use its own large address space.

Two ideas are worth separating at the start:
\begin{itemize}
  \item Virtual memory extends \textbf{capacity}. It lets a program use more
    memory than the physical RAM provides.
  \item Virtual memory does \textbf{not} add speed. Because it relies on the
    disk, it is \textbf{much slower} than RAM.
\end{itemize}

\section{Virtual addresses, physical addresses, and the MMU}
A program does not work with the real location of its data in RAM. It works
with \textbf{virtual addresses}, which form a large, uniform address space.
The hardware translates each virtual address into the actual \textbf{physical
address} in RAM when the instruction executes.

\begin{definitionbox}{Virtual address}
The address generated by a running program; it belongs to the program's large
virtual address space and is not the real RAM location.
\end{definitionbox}

\begin{definitionbox}{Physical address}
The actual location of data in main memory (RAM), which the hardware uses after
translation.
\end{definitionbox}

The translator is the \textbf{MMU}, the Memory Management Unit.

\begin{definitionbox}{MMU (Memory Management Unit)}
The hardware that translates virtual addresses into physical addresses at run
time, using the page table maintained by the operating system.
\end{definitionbox}

The MMU is central to every memory access. It does \emph{not} store files and it
does not itself create storage; it maps one address form to another. Keeping
this role straight avoids two common exam errors: attributing file storage to
the MMU, or describing translation as a software-only operation with no
hardware component.

\begin{example}
A question asks what the MMU does. The correct answer is that it translates
virtual addresses into physical addresses. An option saying that the MMU stores
files permanently is wrong: file storage is the job of secondary storage, not
the address translator. An option saying the MMU holds the BIOS is also wrong;
that describes firmware, not the unit that maps addresses.
\end{example}

\section{Paging: fixed-size pages and frames}
\textbf{Paging} divides memory into equal, fixed-size blocks. The virtual
address space of a process is split into \textbf{pages}, and physical memory is
split into \textbf{frames} of exactly the same size. Any page can be placed in
any free frame, and the page table records which page currently sits in which
frame.

\begin{definitionbox}{Page}
A fixed-size block of the virtual address space in paging.
\end{definitionbox}

\begin{definitionbox}{Frame}
A fixed-size block of physical memory (RAM) into which a page is loaded; frames
are the same size as pages.
\end{definitionbox}

Because every block is the same size, paging makes memory allocation simple and
predictable. There is no awkward leftover space caused by blocks of many
different sizes. The cost is a small amount of unused space in the last page of
a program, but the mechanism itself is uniform.

Two related terms must not be swapped. \textbf{Pages} are the \emph{virtual}
blocks; \textbf{frames} are the \emph{physical} blocks. They are the same size,
but a page lives in the program's virtual view while a frame lives in RAM.

\section{Segmentation: variable-size logical segments}
\textbf{Segmentation} divides a program into \textbf{variable-size segments},
each corresponding to a logical unit of the program: the code, the data, the
stack, a procedure, or a module. Because these logical units differ in size,
segment sizes differ as well.

\begin{definitionbox}{Segment}
A variable-size block of a program that corresponds to a logical unit such as
code, data, stack, or a module.
\end{definitionbox}

Segmentation reflects how a programmer thinks about a program: one segment per
logical part. Paging, by contrast, reflects no logical unit at all --- its pages
are simply equal slices chosen for uniform allocation. This is the key
near-neighbour distinction between the two techniques.

\begin{example}
A stem asks which technique divides a program according to logical units such as
code and data. The answer is segmentation, because those units are logical and
of different sizes. If the stem instead describes equal fixed-size blocks and
the terms ``page'' and ``frame'', the answer is paging.
\end{example}

\section{Paging versus segmentation at a glance}
\begin{center}
\begin{tabular}{p{0.30\textwidth}p{0.30\textwidth}p{0.30\textwidth}}
\toprule
\textbf{Feature} & \textbf{Paging} & \textbf{Segmentation}\\
\midrule
Block size & Fixed & Variable\\
Blocks & Pages (virtual) and frames (physical) & Logical segments\\
Sizes equal? & Yes, all equal & No, they differ\\
Logical view & Not visible to the programmer & Matches program structure\\
Unit & Equal slices & Code, data, stack, module\\
\bottomrule
\end{tabular}
\end{center}

The single most common error is to call paging's blocks ``segments'' or to say
that segmentation uses equal-sized blocks. Keep the pairing firm:
\textit{paging is fixed-size; segmentation is variable-size and logical}.

\section{Page faults and demand paging}
When a program references a page that is \emph{not} currently in RAM, the
hardware raises a \textbf{page fault}. The operating system then loads the
required page from the disk into a free frame and lets the program continue.

\begin{definitionbox}{Page fault}
An event that occurs when a referenced page is not currently in RAM; the page
must be loaded from secondary storage (disk) into a free frame.
\end{definitionbox}

\begin{definitionbox}{Demand paging}
A policy in which a page is loaded into memory only when it is actually
referenced, rather than loading the whole program in advance.
\end{definitionbox}

A page fault is a normal, expected event --- \textbf{not a hardware fault}. It
simply means the needed page is on the disk for the moment. Demand paging is
what makes a program able to start before every one of its pages is in RAM: the
remaining pages are brought in on demand, as execution reaches them.

\begin{example}
A question says a running program references a page that is not in RAM and asks
what happens. The correct answer is a page fault followed by loading that page
from the disk. An option describing a physical hardware failure is wrong,
because a page fault is a routine memory-management event, not a broken
component.
\end{example}

\section{Swapping and the disk-backed hierarchy}
\textbf{Swapping} is the movement of pages or whole processes between RAM and
secondary storage (the disk). Because virtual memory depends on this movement,
it is \textbf{disk-backed}: the disk supplies the extra capacity that RAM does
not have.

\begin{definitionbox}{Swapping}
Moving pages or processes between main memory (RAM) and secondary storage
(disk) so that the memory actually needed at a moment can reside in RAM.
\end{definitionbox}

The memory hierarchy, from fastest and smallest to slowest and largest, runs
from CPU registers, to cache, to RAM, to disk-backed virtual memory. Two facts
follow directly:
\begin{enumerate}
  \item Virtual memory is \textbf{much slower than RAM}, because the disk is far
    slower than main memory.
  \item Virtual memory is \textbf{not additional physical RAM}. It borrows disk
    space; it does not add RAM chips.
\end{enumerate}

This is also where virtual memory is most often confused with \textbf{cache}.
They are opposite in purpose and location. Cache is small, fast memory sitting
close to the CPU and it \emph{speeds up} repeated access to hot data. Virtual
memory is large, disk-backed, and it \emph{extends capacity} using slower
storage. Cache helps speed; virtual memory helps size.

\begin{example}
A stem contrasts cache and virtual memory. The correct statement is that cache
is small, fast memory close to the CPU, while virtual memory is slower and
disk-backed. Options that call both ``secondary storage'', or that say virtual
memory speeds up the CPU, reverse the roles and are wrong.
\end{example}

\section{Thrashing}
\textbf{Thrashing} is \emph{excessive paging}. When too many processes compete
for too little RAM, the system spends most of its time moving pages between RAM
and disk, and very little time doing useful work. The CPU stays busy, yet
throughput collapses.

\begin{definitionbox}{Thrashing}
A condition of excessive paging in which the system spends most of its time
swapping pages between RAM and disk, leaving little time for useful work.
\end{definitionbox}

Thrashing is a symptom of memory pressure, not a hardware failure. Adding
physical RAM, or reducing the number of active processes, usually relieves it.

\section{Exam retrieval and common confusions}
Use two retrieval directions. Given the technique, recall its block rule; given
a description, name the technique.
\begin{itemize}
  \item Virtual memory $\leftrightarrow$ extends memory using disk; slower than
    RAM; not extra physical RAM.
  \item MMU $\leftrightarrow$ Memory Management Unit; translates virtual to
    physical addresses.
  \item Paging $\leftrightarrow$ fixed-size pages (virtual) and frames
    (physical).
  \item Segmentation $\leftrightarrow$ variable-size logical segments.
  \item Page fault $\leftrightarrow$ referenced page not in RAM; load it from
    disk.
  \item Demand paging $\leftrightarrow$ load a page only when it is referenced.
  \item Swapping $\leftrightarrow$ move pages/processes between RAM and disk.
  \item Thrashing $\leftrightarrow$ excessive paging.
\end{itemize}

Three distinctions prevent most errors in this topic:
\begin{enumerate}
  \item Do not say virtual memory is \textbf{extra physical RAM}. It is
    disk-backed and \textbf{slower} than RAM; it adds capacity, not speed.
  \item Do not confuse \textbf{paging} (fixed-size pages and frames) with
    \textbf{segmentation} (variable-size logical segments).
  \item Do not call a \textbf{page fault} a hardware fault. It means the
    referenced page must be loaded from the disk; and do not confuse
    \textbf{cache} (small, fast, near the CPU) with virtual memory (large,
    disk-backed).
\end{enumerate}

\subsection*{Exercises}
\begin{enumerate}[label=\arabic*.]
  \item Explain why virtual memory is said to extend memory but not to add
    physical RAM.
  \item Expand MMU and state its function in one sentence.
  \item Distinguish paging from segmentation in terms of block size and the
    unit being divided.
  \item What is a page fault, and what happens after one occurs?
  \item Define thrashing and give one practical way to reduce it.
\end{enumerate}

\subsection*{Solutions}
\begingroup\small
\begin{enumerate}[label=\arabic*.]
  \item Virtual memory extends the addressable memory by using secondary
    storage (disk) to hold parts of a program not currently in RAM. It does not
    add RAM chips, so it is not additional physical RAM; because it is
    disk-backed, it is much slower than RAM.
  \item MMU is the Memory Management Unit. It translates virtual addresses
    generated by a program into physical addresses in RAM at run time.
  \item Paging uses fixed-size, equal blocks: pages in the virtual space and
    frames in physical memory. Segmentation uses variable-size blocks that
    correspond to logical program units such as code, data, or a module.
  \item A page fault occurs when a program references a page that is not
    currently in RAM. The operating system then loads that page from the disk
    into a free frame and resumes execution.
  \item Thrashing is excessive paging, in which the system spends most of its
    time swapping pages between RAM and disk and little time on useful work.
    Adding physical RAM (or reducing the number of active processes) usually
    reduces it.
\end{enumerate}
\endgroup
\end{document}
